\section{Часть 2 — Метод Baldur}

\begin{frame}{Что такое LLM и почему это важно}
\textbf{LLM} — большая языковая модель (Transformer), обученная на больших корпусах текста/кода:
\begin{itemize}
  \item генерирует текст и код;
  \item умеет «достраивать» последовательности;
  \item может быть дообучена на специализированных данных (fine-tuning).
\end{itemize}

\vspace{2mm}
В статье используется \textbf{Minerva} — модель, предварительно обученная на математических текстах.
\end{frame}

\begin{frame}{Whole-proof generation: идея}
Вместо пошагового поиска:
\begin{itemize}
  \item вход: формулировка теоремы;
  \item выход: \textbf{полный proof script}.
\end{itemize}

\vspace{2mm}
Дальше proof assistant просто проверяет:
\begin{itemize}
  \item нет ошибок $\Rightarrow$ теорема доказана;
  \item есть ошибки $\Rightarrow$ пробуем другой сэмпл или включаем repair.
\end{itemize}
\end{frame}


\begin{figure}{}
\centering
\input{figures/figure1}
\caption{An example of using the proof generation model to generate a proof.}
\end{figure}

% \begin{frame}[fragile]{Псевдокод: генерация нескольких попыток}
% \lstset{style=code}
% \begin{lstlisting}[language=Python]
% for temp in [0.0, 0.4, 0.8, 1.0]:
%     proofs = sample_llm(theorem, k=64, temperature=temp, top_k=40)
%     for p in deduplicate(proofs):
%         ok = isabelle_check(p)
%         if ok:
%             return SUCCESS, p
% return FAILURE
% \end{lstlisting}
% \end{frame}

\begin{frame}{Proof repair: как Baldur чинит доказательства}
Если Isabelle возвращает ошибку, Baldur формирует вход для repair-модели:
\begin{itemize}
  \item теорема;
  \item неверный proof script;
  \item \textbf{сообщение об ошибке} (error message).
\end{itemize}

\vspace{2mm}
Далее модель генерирует исправленный proof script и снова запускается проверка.
\end{frame}

\begin{frame}[fragile]{Формат repair-входа (пример)}
\lstset{style=code, language=}
\begin{lstlisting}
<THEOREM>
lemma fun_sum_commute:
  assumes "f 0 = 0" and "forall x y. f (x + y) = f x + f y"
  shows "f (sum g A) = (\sum a \in A. f (g a))"

<INCORRECT_PROOF>
proof (induct A)
  ...
qed

<ERROR>
Step error: Unable to figure out induct rule
At command "proof" (line 1)
\end{lstlisting}
\end{frame}

\begin{frame}{Контекст из файла теории (theory file context)}
Дополнительный сигнал: строки \textbf{перед теоремой} в том же файле.
Они могут содержать:
\begin{itemize}
  \item определения;
  \item предыдущие теоремы;
  \item доказательства;
  \item комментарии.
\end{itemize}

\vspace{2mm}
Это помогает модели копировать/адаптировать похожие конструкции рядом.
\end{frame}

\begin{frame}[fragile]{Пример влияния контекста}
\lstset{style=code, language=}
\begin{lstlisting}
lemma additive_implies_homogenous:
  assumes "forall x y. f (x + y) = f x + f y"
  shows   "f 0 = 0"
by (metis add.left_neutral assms)

lemma fun_sum_commute:
  assumes "f 0 = 0" and "forall x y. f (x + y) = f x + f y"
  shows "f (sum g A) = (\sum a \in A. f (g a))"
\end{lstlisting}
\end{frame}

\begin{frame}{Почему подход проще, чем step-by-step search}
Step-by-step:
\begin{itemize}
  \item много обращений к модели;
  \item постоянный цикл «модель → prover → модель → ...»;
  \item сложнее инфраструктура.
\end{itemize}

\vspace{2mm}
Whole-proof:
\begin{itemize}
  \item один запрос = много шагов сразу;
  \item удобно батчить генерацию;
  \item проверка отдельно на CPU.
\end{itemize}
\end{frame}